Audit kepatuhan terhadap ISO/IEC 27001:2022 memerlukan pemetaan dokumen kebijakan organisasi terhadap kontrol spesifik pada Annex A, suatu proses yang memakan waktu dan rentan inkonsistensi apabila dilakukan secara manual. Penelitian ini mengembangkan dan mengevaluasi strategi \textit{information retrieval} dalam kerangka kerja \textit{Retrieval-Augmented Generation} (RAG) untuk mendukung pemetaan kontrol \textit{Statement of Applicability} (SoA) secara otomatis pada ISO/IEC 27001:2022. Tiga sumbu desain dievaluasi secara sistematis: (1) komposisi bobot antara \textit{dense retrieval} dan \textit{sparse retrieval} pada \textit{hybrid retrieval}, (2) ukuran \textit{candidate pool} untuk \textit{reranking cross-encoder}, dan (3) bobot \textit{score fusion} antara skor \textit{hybrid retrieval} dan skor \textit{cross-encoder}.

Eksperimen dilakukan menggunakan 77 \textit{query} yang diturunkan dari empat dokumen kebijakan Direktorat Teknologi Informasi (DTI) Universitas Gadjah Mada terhadap basis pengetahuan 93 kontrol Annex A (A.5 sampai A.8) dari ISO/IEC 27001:2022. Basis pengetahuan terstruktur sebagai entitas JSON yang memuat atribut deskriptif, indikator audit, kata kunci, aset terkait, dan prinsip keamanan dalam bahasa Inggris dan Indonesia. Komponen \textit{sparse retrieval} menggunakan BM25, komponen \textit{dense retrieval} menggunakan model \textit{embedding} paraphrase-multilingual-MiniLM-L12-v2, dan komponen \textit{reranking} menggunakan \textit{cross-encoder} BAAI/bge-reranker-v2-m3. Metrik evaluasi mencakup Hit@3, Recall@3, \textit{Mean Average Precision} (MAP@3), dan NDCG@3, dilengkapi uji statistik non-parametrik Friedman.

Hasil eksperimen menunjukkan bahwa konfigurasi \textit{hybrid retrieval} seimbang dengan \textit{dense}=0{,}50 dan \textit{sparse}=0{,}50 menghasilkan performa terbaik pada tiga dari empat metrik (Recall@3=0{,}537; MAP@3=0{,}467; NDCG@3=0{,}575). Ukuran \textit{candidate pool} sebesar 5 memberikan kualitas \textit{reranking} tertinggi dengan penurunan performa monoton seiring bertambahnya ukuran \textit{pool}. Komposisi \textit{score fusion} dominan \textit{hybrid} dengan Hybrid:CE=0{,}75:0{,}25 memberikan performa keseluruhan terbaik (Recall@3=0{,}541; NDCG@3=0{,}579). Uji Friedman mengonfirmasi bahwa perbedaan performa antar konfigurasi signifikan secara statistik ($p < 0{,}05$) pada seluruh skenario eksperimen. Analisis per-\textit{query} mengungkapkan bahwa \textit{hybrid retrieval} dan \textit{reranking cross-encoder} memiliki kekuatan komplementer: \textit{hybrid retrieval} unggul pada \textit{query} yang mengandung terminologi teknis eksplisit (36 dari 77 \textit{query}), sementara \textit{cross-encoder} menguntungkan \textit{query} yang memerlukan pemahaman kontekstual lebih mendalam (15 dari 77 \textit{query}).

Penelitian ini menyimpulkan bahwa konfigurasi optimal dari \textit{hybrid retrieval} seimbang, \textit{candidate pool} berukuran kecil, dan \textit{score fusion} dominan \textit{hybrid} mencapai keseimbangan terbaik antara kualitas \textit{retrieval} dan efisiensi komputasi untuk pemetaan kontrol SoA ISO/IEC 27001:2022.

\noindent{Kata kunci} : ISO/IEC 27001, \textit{Statement of Applicability}, \textit{hybrid retrieval}, \textit{cross-encoder reranking}, \textit{retrieval-augmented generation}, audit keamanan informasi
